\section{Canonical quotients and an elementary valuation bound}
We refine the residual comparisons by retaining quotient markings, then obtain an elementary proof of the supersingular upper step at \(p=5\) and \(e=20\).

\subsection{Rational points and quotient markings}
This subsection records the rationality and marking facts needed for the two comparisons.
Keep the notation of Lemma~\ref{rc-lem:newton}: \(p\geq5\), \(e=v_K(p)\), \(L/K\) is the common tame extension of degree \(m=p^2-1\), \(R=\Ocal_L\), and \(E=me\). The residue field is algebraically closed. In the two-slope range put \(h=v_L(u)\), \(A=(E-h)/(p-1)\), and \(B=h/[p(p-1)]\). Choose \(\pi,\tau,\zeta\) with \(\tau(\pi)=\zeta\pi\) as before.

\begin{lemma}\label{fc-lem:rational}
The canonical \(p\)-torsion points in the two-slope case are rational over
\(L\). All \(p\)-torsion points in the one-slope case are rational over
\(L\). In the two-slope case, if a noncanonical \(p\)-line is also
\(L\)-stable, then all \(p\)-torsion is rational over \(L\).
\end{lemma}
\begin{proof}
The high valuation \(A=(E-h)/(p-1)\) is integral by \eqref{rc-eq:hasseweight}.
Substituting \(T=\pi^A X\) in \([p](T)/T\) and dividing by its smallest
coefficient valuation gives a nonzero constant plus a unit times
\(X^{p-1}\) on reduction. Its nonzero roots are simple and lie in \(k\),
so Hensel's lemma gives the canonical points over \(L\).

In the one-slope case substitute \(T=\pi^eX\) in \([p](T)\) and divide
by \(\pi^{p^2e}\). The reduction is an additive polynomial of degree
\(p^2\), with nonzero linear coefficient given by the residue of
\(p/\pi^E\). The coefficient of \(X^p\) may vanish, and that of
\(X^{p^2}\) is a unit. The terms contributed by formal addition have
larger valuation. Thus this polynomial has \(p^2\) distinct roots in
\(k\), which all lift over \(L\). This includes equality in the
one-slope boundary.

Finally, the determinant of the good height-two, dimension-one Tate
module over \(L\) is cyclotomic. Since \(\mu_p\subset L\), its reduction
modulo \(p\) is trivial. The canonical line is already constant, so its
quotient character is trivial. A second stable line splits these two
trivial characters, proving the last assertion.
\end{proof}


For a canonical quotient \(q:G\to G/C\), the dual kernel \(q(G[p])\) is noncanonical when \(h<E/(p+1)\), canonical when \(h>E/(p+1)\), and the quotient is one-slope at equality. Indeed a low \(p\)-point has quotient valuation \(h/(p-1)\); comparison with the quotient slopes in Proposition~\ref{rc-prop:quotients} gives exactly these alternatives. That proposition also proves that every \(p\)-quotient of a one-slope group has Hasse valuation \(E/(p+1)\).

An isomorphism of \(p^2\)-torsion carrying \(C\) to \(C'\) induces an isomorphism of the quotient \(p\)-torsion including these marked dual lines. The quotient groups are the subquotients
\[
\{x\in G[p^2]:px\in C\}/C,\qquad
\{y\in H[p^2]:py\in C'\}/C',
\]
and the markings are the images of \(G[p]\) and \(H[p]\).

\subsection{Two canonical-neighbor comparisons}
These comparisons retain the dual-kernel marking when passing through a canonical quotient.


\begin{proposition}\label{fc-prop:central}
Suppose \(p\mid E\) and \(G[p^2]\simeq H[p^2]\) over \(K\).
If \(G\) has a two-slope Hasse valuation \(h_G\) prime to \(p\) with
\(h_G>E/(p+1)\), then the special inertia traces agree.
\end{proposition}
\begin{proof}
The residual splitting field of \(G\) is nontrivial. Therefore \(H\) is
two-slope. If \(p\nmid h_H\), the residual comparison of Section~\ref{rc-sec:residual} gives the conclusion.
Otherwise \(p\mid h_H\) and \(h_H<h_G\). Wild inertia identifies the
canonical lines. Quotient by them to obtain \(G_1,H_1\), whose residual
representations and marked dual lines are identified.

We have \(h_{G_1}=E-h_G\), prime to \(p\), and its dual marking is
canonical. If \(h_H>E/(p+1)\), then \(h_{H_1}=E-h_H\) is divisible by
\(p\) and greater than \(h_{G_1}\), contradicting
the residual comparison of Section~\ref{rc-sec:residual}. If \(h_H<E/(p+1)\), the marking on \(H_1\) is
noncanonical, whereas it must correspond to the unique wild-fixed
canonical line of \(G_1\). Equivalently, Lemma~\ref{fc-lem:rational} makes
\(H_1[p]\) trivial over \(L\), again a contradiction. At
\(h_H=E/(p+1)\), the quotient \(H_1\) is one-slope and has trivial
residual representation over \(L\), with the same contradiction.
\end{proof}

\begin{proposition}\label{fc-prop:oneslope}
Suppose \(p\nmid E\) and \(G[p^2]\simeq H[p^2]\) over \(L\). If \(G\)
is one-slope, then \(H\) is one-slope.
\end{proposition}
\begin{proof}
Both groups have all \(p\)-torsion rational over \(L\). If \(H\) were
two-slope, its low-point valuation would force \(p\mid h_H\).
Transport its canonical line to a line of \(G[p]\) and quotient both.
Every quotient on the \(G\) side has Hasse valuation \(E/(p+1)\), prime
to \(p\), and residual lower break exactly \(E\).

For \(h_H<E/(p+1)\), the marked dual line of the canonical quotient of
\(H\) is noncanonical and stable. Its residual representation is trivial,
which is impossible. For \(h_H>E/(p+1)\), the quotient Hasse valuation
\(E-h_H\) is prime to \(p\); its break is
\[
 \frac{(p+1)h_H-E}{p-1}<E,
\]
also impossible. Equality \(h_H=E/(p+1)\) is incompatible with
\(p\mid h_H\). Thus \(H\) is one-slope.
\end{proof}

\begin{remark}
Proposition~\ref{fc-prop:oneslope} recognizes a valuation regime over the
common good field. It does not, by itself, recover descent characters
over \(K\).
\end{remark}

\begin{corollary}\label{fc-cor:six}
Let \(p=5\) and \(e=20\). Consider a hypothetical unequal-type pair with
isomorphic \(p^2\)-torsion, one two-slope Hasse valuation prime to \(5\) and
the other divisible by \(5\). Write \(h_i=4z_i\), with \(z_G\) prime to \(5\).
Then its possible ordered normalized valuations are confined to
\[
 (z_G,z_H)\in
 \{(9,5),(13,5),(17,5),(14,10),(18,10),(19,15)\}.
\]
\end{corollary}
\begin{proof}
Here \(E=480\), the two-slope range is \(1\leq z\leq99\), and the
outer range is \(z<20\). the residual comparison of Section~\ref{rc-sec:residual} and
Proposition~\ref{fc-prop:central} leave only
\(0<z_H<z_G<20\), with \(5\mid z_H\) and \(5\nmid z_G\).
The common canonical character in \eqref{rc-eq:kappa} gives
\(a_H-a_G\equiv z_G-z_H\pmod{24}\).
The finite determinant characters agree, so
\(\alpha_G^6=\alpha_H^6\); hence \(4\mid z_G-z_H\).
For \(z_H=5,10,15\), the listed possibilities follow immediately.
No occurrence of a listed pair at any level is asserted.
\end{proof}

\subsection{The graph input and its compatible flag}

We use the graph-specialization statement proved in the proof of Theorem~\ref{thm:common-determinant}.
If an actual isomorphism \(G[p^n]_L\simeq H[p^n]_L\) descends to \(K\)
and the special inertia traces differ, its finite flat graph closure
\(X_n\subset G[p^n]\times H[p^n]\) satisfies
\begin{equation}\label{fc-eq:graphinput}
 (X_n)_k=\ker(F_G^n)\times\ker(F_H^n).
\end{equation}
Here \(F^n\) is the iterated relative Frobenius of the special formal
group. The proof of that statement, based on disjoint descent-character
orbits and maximal isotropy of the graph, is not repeated here.

\begin{proposition}\label{fc-prop:flag}
For \(0\leq i\leq n\), let \(X_i\) be the closure of the restricted
graph at level \(p^i\), with \(X_0=0\). Then these form a finite flat
flag, with \((X_i)_k=\ker(F_G^i)\times\ker(F_H^i)\).
The quotients \(Q_i=X_i/X_{i-1}\) have rank \(p^2\), are killed by \(p\),
and specialize to \(\alpha_p^2\). For \(i\geq2\), multiplication by \(p\)
induces \(Q_i\to Q_{i-1}\) which is a generic-fiber isomorphism and is
zero on the special fiber.
\end{proposition}
\begin{proof}
The generic graph inclusions extend to their flat closures. Multiplication
by \(p\) maps \(X_i\) to \(X_{i-1}\), since the factorization holds
generically and the source is flat. The quotients are finite flat of rank
\(p^2\) and are killed by \(p\), again by generic-fiber faithfulness.
Apply \eqref{fc-eq:graphinput} at each level. Its successive special
quotients are the first Frobenius kernels of the corresponding Frobenius
twists of the two supersingular groups. Each is \(\alpha_p\): it and its
dual are connected, excluding the etale and multiplicative order-\(p\)
groups. The generic multiplication maps on successive quotients of
\((\mathbf Z/p^i)^2\) are isomorphisms. Specially, \([p]\) has initial
degree \(p^2\), so it sends \(\ker(F^i)\) into \(\ker(F^{i-2})\) and
therefore induces zero on the asserted successive quotients.
\end{proof}

\subsection{An invariant augmentation-ideal bound}
The graph algebra gives an elementary numerical obstruction without choosing a canonical subgroup.


\begin{proposition}\label{fc-thm:valuation}
Let \(M,H/K\) be potentially supersingular height-two Barsotti--Tate
groups with isomorphic \(p^n\)-torsion and unequal special inertia traces.
Over the common tame good field $L$ of degree $p^2-1$, every nonzero $p^n$-torsion coordinate in either good integral formal model satisfies
\[
 v_L(T(P))\ge\frac{p-1}{p^n-1}.
\]
Consequently,
\begin{equation}\label{fc-eq:mainbound}
 e(p^n-1)\geq(p-1)p^{2n-2}.
\end{equation}
No assumption on \(v_p(e)\) is required.
\end{proposition}
\begin{proof}
Use the common good field \(L\) and the notation above. The determinant
comparison gives
\(\alpha^{p+1}=\beta^{p+1}\) for the two tangent descent characters.
Indeed both determinant Tate characters have the same cyclotomic factor,
and their remaining finite inertia factors have order prime to \(p\).
The torsion isomorphism makes these factors congruent modulo \(p^n\),
hence equal. The good dimension-one determinant identity is recorded in
the proof of \citet[Proposition~5.2.2]{Liu2007}. Inverting both character
conventions leaves the displayed equality unchanged.

Let \(\Delta\subset\Gal(L/K)\) have order \(p-1\), generated by
\(\delta=\tau^{p+1}\). Its two special tangent scalars are the same
\(c\in\mu_{p-1}\subset\Z_p^*\). Put \(R_0=R^\Delta\). Every nonzero
element of its maximal ideal has \(v_L\) at least \(p-1\).

We first choose exact scalar coordinates. For a formal parameter \(T\),
the projector
\[
 \frac1{p-1}\sum_{j=0}^{p-2}c^{-j}\delta^j(T)
\]
has unit linear coefficient and is a \(c\)-eigenparameter. The canonical
scalar multiplications \([a]\), for \(a\in\mu_{p-1}\), commute with
descent. Applying also
\[
 \frac1{p-1}\sum_{a\in\mu_{p-1}}a^{-1}[a]^*(T)
\]
preserves that property and makes \([a]^*T=aT\). Both projectors preserve
a unit tangent term. Choose such parameters \(T,S\) for the two formal
groups. Then
\[
 \delta(T)=cT,\quad\delta(S)=cS,\qquad
 [a]^*T=aT,\quad[a]^*S=aS.
\]
If pullback conventions give the simultaneous scalar \(c^{-1}\), the same
construction uses that scalar instead. Integral parameter changes with
unit linear term preserve all positive point valuations.

Let \(A_X\) be the algebra of the graph closure, with augmentation ideal
\(I_X\), and write \(x,y\) for the coordinate images. By
\eqref{fc-eq:graphinput},
\[
 A_X/\pi A_X=k[x,y]/(x^{p^n},y^{p^n}).
\]
Thus the monomials \(x^iy^j\), for \(0\leq i,j<p^n\), form an
\(R\)-basis: they generate by Nakayama and their number equals the
finite free rank \(p^{2n}\). Express uniquely
\[
 x^{p^n}=\sum c_{ij}x^iy^j,
 \qquad y^{p^n}=\sum d_{ij}x^iy^j.
\]
Every coefficient is in the maximal ideal of \(R\), and
\(c_{00}=d_{00}=0\) by augmentation.

The graph is stable under descent and simultaneous scalar multiplication.
The semilinear operator \([c]^{*\,-1}\circ\delta\) therefore preserves
\(A_X\), acts as \(\delta\) on \(R\), and fixes \(x,y\). Uniqueness of
the two basis expansions forces \(c_{ij},d_{ij}\in R_0\). Consequently
\begin{equation}\label{fc-eq:augmentation}
 x^{p^n},y^{p^n}\in\pi^{p-1}I_X.
\end{equation}

For a nonzero torsion point \(P\) and its nonzero graph image, put
\[
 r=\min\{v_L(T(P)),v_L(S(f(P)))\}>0.
\]
Every nonconstant monomial has
valuation at least \(r\). Evaluating the two expansions in
\eqref{fc-eq:augmentation} gives
\[
 p^n r\geq p-1+r,\qquad
 r\geq\frac{p-1}{p^n-1}.
\]
Cancellation can only increase the valuation of either sum. Thus every
nonzero \(p^n\)-torsion coordinate of both groups satisfies this bound.

Finally consider
\[
 Q_n(T)=\left.\frac{[p](U)}U\right|_{U=[p^{n-1}](T)}.
\]
Its constant term is \(p\), and it is distinguished of degree
\(D=(p^2-1)p^{2n-2}\), because \([p]\) has special initial degree
\(p^2\). Its Weierstrass roots are precisely the primitive \(p^n\)-torsion
coordinates. The constant term of its Weierstrass polynomial has
valuation \(E\), so the sum of their valuations is \(E\). Their average
is therefore
\[
 \frac ED=\frac e{p^{2n-2}}.
\]
Comparing this average with the lower bound for every root proves
\eqref{fc-eq:mainbound}.
\end{proof}

\begin{corollary}\label{fc-cor:Aupper}
For \(p=5\) and \(e_K=20\), two potentially supersingular height-two
groups with isomorphic \(125\)-torsion have equal special inertia traces.
\end{corollary}
\begin{proof}
Otherwise Proposition~\ref{fc-thm:valuation}, with \(n=3\), would give
\[
 20(125-1)\geq4\cdot5^4,
 \qquad\text{that is,}\qquad2480\geq2500.
\]
This is false. Equivalently, the primitive-torsion average is \(4/125\),
whereas every root would have to have valuation at least \(4/124\).
\end{proof}

Together with the already-proved structural and mixed comparisons and the
unequal-trace \(25\)-torsion construction, Corollary~\ref{fc-cor:Aupper} gives
\(\nu_{\min}(5,20,2)=2\) in the convention where index \(\nu\) compares
level \(p^{\nu+1}\). The pointwise estimate in Proposition~\ref{fc-thm:valuation} also combines with the central point-field criterion in Section~\ref{sec:central-ramification} to prove the general prime-to-$p$ recovery theorem and its stronger supersingular range. That argument uses actual canonical quotients and transported kernels; it does not require lifting the special graph splitting in \eqref{fc-eq:graphinput} to generic cyclic subgroups.
